\documentclass[a4paper,11pt]{article}
\usepackage[brazil]{babel}
\usepackage[utf8]{inputenc}
\usepackage[T1]{fontenc}
\usepackage{amsmath,amssymb,amsthm}
\usepackage{geometry}
\usepackage{hyperref}
\usepackage{booktabs}
\usepackage{xcolor}
\usepackage{tcolorbox}
\usepackage{enumitem}
\usepackage{microtype}
\usepackage{parskip}
\usepackage{titlesec}
\usepackage{fancyhdr}
\usepackage{mdframed}
\usepackage{mathtools}
\usepackage{algorithm}
\usepackage{algpseudocode}

\geometry{margin=2.5cm}

\definecolor{azul}{RGB}{20,60,120}
\definecolor{cinza}{RGB}{245,245,245}
\definecolor{verde}{RGB}{30,130,70}
\definecolor{vermelho}{RGB}{160,30,30}
\definecolor{laranja}{RGB}{180,100,0}

\pagestyle{fancy}
\fancyhf{}
\fancyhead[L]{\small\textcolor{gray}{Teste Estatístico de Alinhamento Semântico}}
\fancyhead[R]{\small\textcolor{gray}{SLO e a Coloração Grundy}}
\fancyfoot[C]{\small\thepage}
\renewcommand{\headrulewidth}{0.4pt}

\titleformat{\section}{\large\bfseries\color{azul}}{}{0em}{}[\titlerule]
\titleformat{\subsection}{\normalsize\bfseries\color{azul!80}}{}{0em}{}

\tcbuselibrary{skins,breakable}
\newtcolorbox{definicao}[1]{
  colback=azul!5, colframe=azul!50,
  fonttitle=\bfseries\small, title=Definição: #1,
  breakable, left=6pt, right=6pt
}
\newtcolorbox{intuicao}{
  colback=verde!5, colframe=verde!50,
  fonttitle=\bfseries\small, title=Intuição,
  breakable, left=6pt, right=6pt
}
\newtcolorbox{atencao}{
  colback=laranja!8, colframe=laranja!60,
  fonttitle=\bfseries\small, title=Atenção,
  breakable, left=6pt, right=6pt
}
\newtcolorbox{resultado}{
  colback=verde!7, colframe=verde!60,
  fonttitle=\bfseries\small, title=Resultado,
  breakable, left=6pt, right=6pt
}

\newtheorem{proposicao}{Proposição}
\newtheorem{observacao}{Observação}

% ------------------------------------------------------------------
\begin{document}

\begin{center}
  {\LARGE\bfseries\color{azul}
    Validação Estatística do Alinhamento Semântico\\[4pt]
    da \textit{Smallest-Last Ordering} (SLO)\\[6pt]
    para a Coloração Grundy}\\[10pt]
  {\normalsize Explicação pedagógica do teste de permutação\\
  e sua relação com a escolha de ordenação na formulação de representantes}\\[4pt]
  {\small \today}
\end{center}

\noindent\rule{\linewidth}{1pt}

% ==================================================================
\section{Contexto: por que a ordenação importa}

Na formulação de programação inteira por representantes assimétricos
para a Coloração Grundy, o conjunto de variáveis de decisão é indexado
por
\[
  \mathcal{X} = \bigl\{(u,v) \;\big|\; v \in \overline{N}(u),\;
  \sigma(u) \leq \sigma(v)\bigr\},
\]
onde $\sigma : V \to \{0,\ldots,n-1\}$ é uma ordem total sobre os
vértices. Para cada aresta $\{u,v\} \in \overline{E}$, exatamente um
dos pares $(u,v)$ ou $(v,u)$ pertence a $\mathcal{X}$: aquele em que
o primeiro elemento tem menor $\sigma$. O vértice de menor $\sigma$
em cada classe de cor é, portanto, o \emph{representante canônico}
dessa classe no modelo.

A hipótese de alinhamento semântico afirma:

\begin{center}
\itshape
A SLO, ao atribuir índices menores a vértices de maior grau residual,
posiciona como representantes canônicos exatamente os vértices que
tendem a ser representantes nas soluções ótimas da Coloração Grundy.
\end{center}

Para testar essa hipótese de forma rigorosa, utilizamos um
\textbf{teste de permutação}, descrito nas seções seguintes.

% ==================================================================
\section{A estatística de teste: rank médio normalizado}

\subsection{Definição}

Seja $G = (V,E)$ um grafo com $n$ vértices e seja $\text{OPT}(G)$ o
conjunto de todas as colorações Grundy ótimas de $G$, isto é, aquelas
que utilizam $\Gamma(G)$ cores. Para uma ordem $\sigma$ e uma coloração
$C \in \text{OPT}(G)$ com classes $V_1^C, \ldots, V_{\Gamma}^C$,
define-se o \emph{representante canônico} da classe $i$ como
\[
  r_i^C(\sigma) \;=\; \arg\min_{v \in V_i^C}\; \sigma(v).
\]

\begin{definicao}{Rank médio normalizado $\bar{r}(\sigma)$}
O \emph{rank médio normalizado} da ordem $\sigma$ é
\[
  \bar{r}(\sigma) \;=\;
  \frac{1}{|\text{OPT}(G)| \cdot \Gamma(G)}
  \sum_{C \in \text{OPT}(G)} \sum_{i=1}^{\Gamma(G)} \sigma\!\left(r_i^C(\sigma)\right)
  \;-\; \frac{\Gamma(G)-1}{2},
\]
onde o termo $(\Gamma(G)-1)/2$ é a normalização pelo melhor caso
possível: se os $\Gamma(G)$ representantes ocupassem as posições
$0, 1, \ldots, \Gamma(G)-1$, o rank médio bruto seria exatamente
$(\Gamma(G)-1)/2$.
\end{definicao}

\begin{intuicao}
$\bar{r}(\sigma) = 0$ significa que, em média sobre todas as soluções
ótimas, os representantes canônicos ocupam as $\Gamma$ primeiras
posições da ordem $\sigma$ — o melhor cenário possível.
$\bar{r}(\sigma) > 0$ significa que os representantes estão, em média,
deslocados para posições mais tardias.
Valores menores de $\bar{r}(\sigma)$ indicam melhor alinhamento
semântico da ordem $\sigma$ com a estrutura das soluções ótimas.
\end{intuicao}

\subsection{Por que não basta comparar com uma única ordem aleatória}

Uma abordagem ingênua compararia $\bar{r}(\sigma_{\text{SLO}})$ com
$\bar{r}(\sigma_{\text{rnd}})$ para uma única permutação aleatória
$\sigma_{\text{rnd}}$. Essa abordagem é estatisticamente inadequada
por duas razões:

\begin{enumerate}
  \item \textbf{Alta variância:} uma permutação aleatória pode por
    acaso ser melhor ou pior que a SLO sem que isso reflita o
    comportamento médio do acaso.
  \item \textbf{Sem significância:} não é possível afirmar que a
    diferença observada é sistemática ou fruto do acaso com apenas
    um ponto de comparação.
\end{enumerate}

O correto é comparar $\bar{r}(\sigma_{\text{SLO}})$ com a
\textbf{distribuição} de $\bar{r}(\sigma)$ quando $\sigma$ é amostrada
uniformemente ao acaso.

% ==================================================================
\section{O teste de permutação}

\subsection{Estrutura do teste}

O teste de permutação é um método \textbf{não paramétrico}: não assume
nenhuma distribuição paramétrica para a estatística de teste. Sua única
hipótese é que, sob a hipótese nula, todas as permutações de $V$ são
igualmente prováveis — o que é exatamente verdade quando $\sigma$ é
amostrada uniformemente.

\medskip
\noindent O teste é estruturado em cinco passos:

\begin{enumerate}[label=\textbf{Passo \arabic*.}, leftmargin=*, itemsep=6pt]

  \item \textbf{Enumeração das soluções ótimas.}
    Calcula-se $\Gamma(G)$ e enumera-se $\text{OPT}(G)$ por força bruta
    (viável para $n \leq 9$, pois há $n!$ permutações a testar).

  \item \textbf{Cálculo do valor observado.}
    Computa-se $\bar{r}(\sigma_{\text{SLO}})$, a estatística de teste
    para a SLO. Este é o valor que queremos avaliar.

  \item \textbf{Construção da distribuição de referência.}
    Geram-se $N$ permutações aleatórias $\sigma_1, \ldots, \sigma_N$
    e calcula-se $\bar{r}(\sigma_k)$ para cada uma, obtendo a amostra
    \[
      \mathcal{R} = \{\bar{r}(\sigma_1), \ldots, \bar{r}(\sigma_N)\}.
    \]
    A média e o desvio padrão dessa amostra estimam o comportamento
    médio do acaso:
    \[
      \mu_{\text{rnd}} = \frac{1}{N}\sum_{k=1}^N \bar{r}(\sigma_k),
      \qquad
      \sigma_{\text{rnd}} = \sqrt{\frac{1}{N-1}\sum_{k=1}^N
      \bigl(\bar{r}(\sigma_k) - \mu_{\text{rnd}}\bigr)^2}.
    \]

  \item \textbf{Cálculo do p-valor empírico.}
    O p-valor unilateral à esquerda é a fração de ordens aleatórias
    que obtiveram resultado \emph{tão bom ou melhor} que a SLO:
    \begin{equation}
      \hat{p} \;=\;
      \frac{\#\{k \in \{1,\ldots,N\} : \bar{r}(\sigma_k)
      \leq \bar{r}(\sigma_{\text{SLO}})\}}{N}.
      \label{eq:pval}
    \end{equation}
    Um p-valor pequeno indica que poucas ordens aleatórias conseguem
    igualar ou superar a SLO — evidência de alinhamento semântico.

  \item \textbf{Decisão estatística.}
    Adota-se o critério convencional:
    \[
      \hat{p} < 0{,}05 \implies \text{rejeita-se } H_0
      \quad (\text{resultado significativo}),
    \]
    onde $H_0$ é a hipótese nula de que a SLO não é melhor que o acaso.

\end{enumerate}

\subsection{Hipóteses do teste}

\begin{align*}
  H_0 &: \bar{r}(\sigma_{\text{SLO}}) \geq \mathbb{E}_\sigma[\bar{r}(\sigma)]
       \quad\text{(SLO não é melhor que o acaso)}\\
  H_1 &: \bar{r}(\sigma_{\text{SLO}}) < \mathbb{E}_\sigma[\bar{r}(\sigma)]
       \quad\text{(SLO posiciona representantes mais cedo)}
\end{align*}

O teste é \textbf{unilateral à esquerda}: interessa apenas saber se a
SLO é \emph{melhor} que o acaso, não se é diferente em qualquer direção.

\subsection{Precisão do p-valor empírico}

O p-valor empírico $\hat{p}$ é uma proporção amostral. Seu erro padrão
é
\[
  \text{EP}(\hat{p}) = \sqrt{\frac{\hat{p}(1-\hat{p})}{N}}.
\]
Para $\hat{p} \approx 0{,}09$ (valor típico nos experimentos) e $N=500$:
\[
  \text{EP} = \sqrt{\frac{0{,}09 \times 0{,}91}{500}} \approx 0{,}013.
\]
Portanto, com $N=500$ amostras, o p-valor estimado tem precisão de
$\pm 0{,}013$ — suficiente para distinguir valores próximos de $0{,}05$.
\textbf{Aumentar $N$ além de 500 melhora a precisão, mas não altera o
p-valor verdadeiro.}

% ==================================================================
\section{O z-score: medida contínua de separação}

Independentemente do limiar de $0{,}05$, a separação entre a SLO e o
acaso pode ser medida continuamente pelo \textbf{z-score}:
\[
  z \;=\; \frac{\mu_{\text{rnd}} - \bar{r}(\sigma_{\text{SLO}})}{\sigma_{\text{rnd}}}.
\]
O z-score mede quantos desvios padrão a SLO está abaixo da média
aleatória. A correspondência com o p-valor (aproximação normal) é:
\[
  p \approx \Phi(-z) = 1 - \Phi(z),
\]
onde $\Phi$ é a função de distribuição acumulada normal padrão.

\begin{center}
\renewcommand{\arraystretch}{1.2}
\begin{tabular}{ccc}
\toprule
$z$ & $p \approx \Phi(-z)$ & Interpretação \\
\midrule
$1{,}282$ & $0{,}10$ & Tendência \\
$1{,}645$ & $0{,}05$ & Significativo ($*$) \\
$2{,}326$ & $0{,}01$ & Altamente significativo ($**$) \\
$3{,}090$ & $0{,}001$ & Muito altamente significativo ($***$) \\
\bottomrule
\end{tabular}
\end{center}

% ==================================================================
\section{Resultados experimentais}

Os experimentos foram conduzidos com grafos aleatórios de Erdős-Rényi
$G(n,p)$ para $n=9$ e cinco densidades, com $N=500$ ordens aleatórias
por grafo. A Tabela~\ref{tab:resultados} apresenta os resultados.

\begin{table}[ht]
\centering
\caption{Resultados do teste de permutação para grafos $G(9, p)$.
$\Gamma$: número de Grundy; $|\text{OPT}|$: número de soluções ótimas;
$\bar{r}_{\text{SLO}}$: rank médio normalizado da SLO;
$\mu_{\text{rnd}}$, $\sigma_{\text{rnd}}$: média e desvio padrão da
distribuição aleatória ($N=500$);
$\hat{p}$: p-valor empírico; $z$: z-score; $\Delta = \mu_{\text{rnd}} -
\bar{r}_{\text{SLO}}$.}
\label{tab:resultados}
\renewcommand{\arraystretch}{1.25}
\begin{tabular}{cccccccccc}
\toprule
$p$ & $\Gamma$ & $|\text{OPT}|$ & $\bar{r}_{\text{SLO}}$ &
$\mu_{\text{rnd}}$ & $\sigma_{\text{rnd}}$ &
$\min$ & $\max$ & $\hat{p}$ & $z$ \\
\midrule
$0{,}10$ & 3 & 181\,440 & $0{,}561$ & $0{,}944$ & $0{,}270$ &
  $0{,}167$ & $1{,}689$ & $0{,}082$ & $1{,}42$ \\
$0{,}25$ & 4 & 4\,158   & $0{,}000$ & $1{,}165$ & $0{,}700$ &
  $0{,}000$ & $3{,}000$ & $0{,}094$ & $1{,}66$ \\
$0{,}50$ & 6 & 325      & $0{,}000$ & $0{,}670$ & $0{,}408$ &
  $0{,}000$ & $1{,}833$ & $0{,}092$ & $1{,}64$ \\
$0{,}75$ & 7 & 4\,320   & $0{,}000$ & $0{,}513$ & $0{,}353$ &
  $0{,}000$ & $1{,}571$ & $0{,}126$ & $1{,}45$ \\
$0{,}90$ & 7 & 80\,640  & $0{,}000$ & $0{,}528$ & $0{,}297$ &
  $0{,}000$ & $1{,}286$ & $0{,}060$ & $1{,}78$ \\
\bottomrule
\end{tabular}
\end{table}

\subsection{Leitura dos resultados}

\paragraph{Nenhum resultado individualmente significativo ao nível 5\%.}
Todos os p-valores ficam entre $0{,}06$ e $0{,}13$, acima do limiar
convencional de $0{,}05$. Isso não implica que a hipótese nula é
verdadeira, mas que a evidência individual por grafo é insuficiente
para rejeitá-la com $n=9$.

\paragraph{O sinal é consistente e não aleatório.}
Três observações apontam para um efeito real:
\begin{itemize}
  \item $\Delta > 0$ em todos os 5 casos: a SLO sempre produz rank
    médio menor que a média do acaso.
  \item Os z-scores estão entre $1{,}42$ e $1{,}78$ — consistentemente
    próximos do limiar de significância.
  \item $\bar{r}_{\text{SLO}} = 0$ em 4 de 5 casos: a SLO atinge
    o melhor caso possível (representantes nas posições mais
    prioritárias) com frequência muito maior do que o acaso.
\end{itemize}

\paragraph{O problema do piso zero.}
Quando $\bar{r}_{\text{SLO}} = 0$, qualquer outra ordem que também
atinja 0 contribui para o p-valor — e o histograma mostra que ordens
aleatórias atingem 0 com probabilidade entre $6\%$ e $13\%$. A SLO
não pode melhorar além de 0, então a separação estatística fica
limitada pelo piso da distribuição, não pela qualidade da SLO.

% ==================================================================
\section{Combinação de testes: Método de Fisher}

Quando múltiplos testes independentes produzem resultados individualmente
não significativos mas consistentemente favoráveis, o
\textbf{método de Fisher} combina os p-valores em um único teste global.

\begin{definicao}{Estatística de Fisher}
Dados $k$ testes independentes com p-valores $\hat{p}_1, \ldots,
\hat{p}_k$, a estatística de Fisher é
\[
  \chi^2_F = -2\sum_{i=1}^{k} \ln(\hat{p}_i),
\]
que segue uma distribuição $\chi^2$ com $2k$ graus de liberdade sob
$H_0$ (independência e uniformidade dos p-valores).
\end{definicao}

Aplicando aos cinco p-valores observados:
\begin{align*}
  \chi^2_F &= -2\bigl(\ln 0{,}082 + \ln 0{,}094 + \ln 0{,}092
               + \ln 0{,}126 + \ln 0{,}060\bigr) \\
           &= -2\bigl(-2{,}501 - 2{,}364 - 2{,}386
               - 2{,}071 - 2{,}813\bigr) \\
           &= -2 \times (-12{,}135) \;=\; 24{,}27.
\end{align*}

Com $2 \times 5 = 10$ graus de liberdade, a tabela $\chi^2$ fornece:
\[
  P(\chi^2_{10} > 24{,}27) \approx 0{,}007.
\]

\begin{resultado}
O método de Fisher combina os cinco testes individuais e produz
$\chi^2_F = 24{,}27$ com $p \approx 0{,}007$ ($< 0{,}01$). Isso
indica que, \textbf{globalmente}, a evidência contra a hipótese nula
é altamente significativa: a SLO posiciona os representantes de forma
sistematicamente melhor que o acaso, mesmo que nenhum grafo individual
produza evidência suficiente.
\end{resultado}

% ==================================================================
\section{Implicações para a formulação de representantes}

A validação estatística do alinhamento semântico tem consequências
diretas para a qualidade do modelo de programação inteira.

\subsection{Qualidade da relaxação LP}

Se a SLO posiciona os representantes naturais com $\sigma$ pequeno,
a variável $x_{u,u}$ correspondente é criada para os vértices que
\emph{deveriam} ser representantes. A relaxação LP tende a atribuir
$x_{u,u}$ fracionário próximo de 1 para esses vértices, pois as
restrições M4 penalizam outros arranjos. Isso produz um limitante dual
de melhor qualidade logo no nó raiz da árvore de branch-and-bound.

Formalmente, quanto mais $\bar{r}(\sigma_{\text{SLO}})$ se aproxima
de 0, mais a estrutura de $\mathcal{X}$ reflete a estrutura ótima do
problema, e menor tende a ser o gap de integralidade na raiz:
\[
  \text{gap}_{\text{raiz}} \;=\;
  \frac{\text{OBJ}_{\text{LP}} - \text{OBJ}_{\text{IP}}}
       {\text{OBJ}_{\text{IP}}}.
\]

\subsection{Redução do espaço de busca}

O número de nós explorados pelo solver depende do gap na raiz e da
simetria residual do modelo. Uma ordem que coloca os representantes
naturais mais cedo reduz a ambiguidade: há menos pares de vértices
que poderiam trocar de papel (representante/membro) sem alterar a
solução. Isso se traduz diretamente em menos branching necessário.

\subsection{Coerência semântica}

A SLO não é apenas uma heurística ad hoc: ela tem justificativa
estrutural. Os vértices de maior grau residual (aqueles com $\sigma$
pequeno na SLO) são os mais prováveis representantes porque:
\begin{itemize}
  \item Um representante $u$ da classe $i$ precisa ter pelo menos
    um vizinho em cada classe $j < i$. Isso exige alto grau em $G$.
  \item Na SLO, vértices de alto grau residual têm $\sigma$ pequeno
    e portanto muitos não-vizinhos à frente na ordem — exatamente
    os candidatos a membros de sua classe.
\end{itemize}

O teste de permutação confirma empiricamente essa coerência:
$\bar{r}(\sigma_{\text{SLO}}) = 0$ em 4 de 5 grafos testados,
e o método de Fisher confirma que esse resultado é globalmente
significativo com $p < 0{,}01$.

% ==================================================================
\section{Limitações e extensões}

\begin{atencao}
As conclusões acima são baseadas em grafos pequenos ($n=9$) devido
à necessidade de enumerar todas as soluções ótimas por força bruta.
Para $n \geq 12$, seria necessário um algoritmo exato de
branch-and-bound para computar $\text{OPT}(G)$, o que tornaria o
estudo computacionalmente viável para instâncias maiores — onde o
efeito da SLO provavelmente seria ainda mais pronunciado, pois o
z-score cresce com $n$ à medida que o espaço de permutações aumenta
fatorialmente.
\end{atencao}

Uma extensão natural seria usar a \textbf{correlação de Spearman}
entre $\sigma(v)$ e a frequência com que $v$ atua como representante
em $\text{OPT}(G)$, definida como
\[
  \rho_s = 1 - \frac{6\sum_v d_v^2}{n(n^2-1)},
\]
onde $d_v = \sigma(v) - \text{rank de frequência de }v$ como
representante. Essa métrica é contínua e não sofre do problema do
piso zero, potencialmente tornando o teste mais poderoso para $n$
pequeno.

\end{document}
